\documentclass[11pt,a4paper]{article}
\usepackage[T1]{fontenc}
\usepackage[utf8]{inputenc}
\usepackage{lmodern,microtype}
\usepackage[margin=23mm,headheight=15pt,headsep=8mm,footskip=12mm]{geometry}
\usepackage[table]{xcolor}
\usepackage{graphicx,booktabs,tabularx,array,ragged2e}
\usepackage{fancyhdr,titlesec,enumitem,amssymb}
\definecolor{maroon}{HTML}{800000}
\definecolor{cream}{HTML}{F7F3ED}
\definecolor{muted}{HTML}{646464}
\definecolor{ready}{HTML}{17623D}
\definecolor{readybg}{HTML}{EDF6F0}
\definecolor{focus}{HTML}{A32932}
\definecolor{focusbg}{HTML}{FCF0F0}
\usepackage[colorlinks=true,linkcolor=maroon,urlcolor=maroon]{hyperref}
\hypersetup{pdftitle={LAlgo: What We Need to Do},pdfauthor={PUP COSC 303 - Group 4}}
\setlength{\parindent}{0pt}
\setlength{\parskip}{5pt}
\emergencystretch=3em
\clubpenalty=10000
\widowpenalty=10000
\titleformat{\section}{\Large\bfseries\color{maroon}}{}{0pt}{}
\titleformat{\subsection}{\normalsize\bfseries\color{maroon}}{}{0pt}{}
\titlespacing*{\section}{0pt}{9pt}{7pt}
\titlespacing*{\subsection}{0pt}{8pt}{4pt}
\setlist{leftmargin=1.4em,itemsep=3pt,topsep=3pt,parsep=0pt}
\newcolumntype{Y}{>{\RaggedRight\arraybackslash}X}
\newcolumntype{L}[1]{>{\RaggedRight\arraybackslash}p{#1}}
\newcommand{\code}[1]{\texttt{#1}}
\newcommand{\good}{\textcolor{ready}{\bfseries READY}}
\newcommand{\todo}{\textcolor{focus}{\bfseries FOCUS}}
\newcommand{\checkitem}[1]{\item[$\square$] #1}
\newcommand{\boxnote}[3]{\par\smallskip\begingroup\setlength{\fboxsep}{8pt}\colorbox{#1}{\parbox{\dimexpr\linewidth-2\fboxsep\relax}{\color{#2}\small #3}}\endgroup\par\smallskip}
\newcommand{\readybox}[1]{\boxnote{readybg}{ready}{#1}}
\newcommand{\focusbox}[1]{\boxnote{focusbg}{focus}{#1}}
\newcommand{\neutralbox}[1]{\boxnote{cream}{black}{#1}}
\newcommand{\profquote}[2]{\neutralbox{\textbf{Ma'am:} ``#1''\par{\footnotesize\color{muted}#2}}}
\pagestyle{fancy}
\fancyhf{}
\fancyfoot[L]{\footnotesize\color{muted}LAlgo team guide}
\fancyfoot[R]{\small\color{maroon}\thepage\,/\,6}
\renewcommand{\headrulewidth}{0pt}
\begin{document}

{\fontsize{23}{28}\selectfont\bfseries\color{maroon}LAlgo: what we need to do\par}
{\small\color{muted}Team guide | Switch and case revision\par}
\vspace{2mm}

\section{1. What matters now}
The lexer and EBNF checker support \code{SWITCH}, \code{CASE}, \code{DEFAULT}, and \code{BREAK}. Run \code{make syntax} to check the switch sample.
This guide condenses the four recordings and the checking rubrics. Short,
timestamped excerpts support the requirements; full transcripts stay in the meeting notes.

\readybox{\good\ means the stated portion exists in the current source or documentation.
The lexer passes \textbf{33/33 local tests}, and the syntax checker passes \textbf{13/13 CLI tests}.}
\focusbox{\todo\ means missing work, partial coverage, or a decision we still need to settle.
Use the syntax checker to demonstrate the implemented grammar.}

\subsection{Language decisions for checking}
\begin{enumerate}
\item \textbf{Use \code{WHILE} for single and nested loops.} The EBNF covers both forms.
\item \textcolor{focus}{\textbf{Keep language features through iteration only.}}
Review future-reserved \code{ALGORITHM}, \code{RETURN}, \code{TRACE}, and \code{INVARIANT};
do not expand them into extra supported constructs for this scope.
\item \textbf{Show \code{SWITCH} during checking.} Use a reference selector, at least one case, one final default, and \code{BREAK;} at the end of every arm.
\end{enumerate}

\profquote{Kung FOR or WHILE, bahala kayo. Pero ang rule niya, dalawa.}{Recording 04, 05:24--05:46. She then names single and nested loops.}

\subsection{Dates stated in the meeting}
{\small\renewcommand{\arraystretch}{1.2}
\begin{tabularx}{\linewidth}{@{}L{26mm}Y@{}}
\toprule
\textbf{Date, 2026} & \textbf{Milestone} \\\midrule
12 November & Submit syntax proposal: grammar, examples, and derivations. Face-to-face by default. \\
19 November & Asynchronous project work. \\
26 November & Lexical checking 2 and syntax checking 1. \\
3 December & Asynchronous project work. \\
10 December & Syntax checking 2, semantics checking 1, final project checking and final exam. \\\bottomrule
\end{tabularx}}

\neutralbox{\textbf{Team priority:} scope decisions $\rightarrow$ grammar and per-rule derivations
$\rightarrow$ parser and diagnostics. The iteration limit concerns language features;
syntax and semantic checking milestones still apply. Submission time/platform and
detailed scoring weights are not supplied in these materials. The semantic guide is now included on page 6.}

\clearpage
\section{2. Lexical analyzer: what already works}
\readybox{\good\ \textbf{Recognition and output are implemented.} The rows below map to the
11 lexical-rubric areas. They describe lexical behavior, not syntax or execution.}

{\small\renewcommand{\arraystretch}{1.21}
\begin{tabularx}{\linewidth}{@{}L{8mm}L{35mm}Y@{}}
\toprule
\textbf{\#} & \textbf{Area} & \textbf{Existing evidence --- \textcolor{ready}{READY}} \\\midrule
1 & Input restriction & Accepts \code{.lal} / \code{.LAL}; rejects other extensions before opening the source. Rejects NUL-containing input. \\
2 & Identifiers & Letter first; then letters, digits, or underscores. Case-sensitive. \\
3 & Keywords & The unified DFA recognizes \code{SWITCH}, \code{CASE}, \code{DEFAULT}, and \code{BREAK} as control-flow keywords. \\
4 & Reserved words & Reserved spellings and separate future-reserved token classes exist. Their final scope needs review below. \\
5 & Constant values & Five categories: integer, decimal, character, string, Boolean. \code{TRUE}/\code{FALSE} count as one category. \\
6 & Noise words & Recognizes \code{SET} and \code{THE}. \\
7 & Comments & Recognizes \code{//} line comments and \code{/* ... */} block comments; retains them in the listing. \\
8a & Arithmetic: 7 & \code{+ - * / \% DIV \textasciicircum}. Assignment \code{=} is counted separately. \\
8b & Boolean: 9 & Logical: \code{AND OR NOT}. Relational: \code{== != < > <= >=}. \\
9 & Delimiters/brackets & Recognizes \code{( ) [ ] \{ \}}, comma, semicolon, colon, and dot. Pairing/nesting checks belong to parsing. \\
10 & Invalid tokens & Emits \code{INVALID} and lexical line/column diagnostics for documented malformed lexemes. \\
11 & Output text file & Writes \code{Symbol Table.txt}: token occurrences, classifications, positions, and lexemes; includes comments/invalid entries. \\\bottomrule
\end{tabularx}}

\profquote{Sa lexical analyzer n'yo, na ang output ay symbol table.}{Recording 04, 00:50--00:57. Retain the lexical output file.}

\subsection{Focus before presenting this as complete}
\begin{itemize}
\item \textbf{Show the keyword tokens and the colon.} Use \code{switch\_sample.lal} and its symbol table during checking.
\item \textcolor{focus}{\textbf{Resolve rubric interpretation.}} The extension filter works, but global
uniqueness of \code{.lal} is unverified. Confirm the intended keyword/reserved distinction
and nine-Boolean-operator classification.
\item \textcolor{focus}{\textbf{Reconcile previous feedback.}} Local tests do not establish that every item
from Ma'am's previous checking has been addressed.
\end{itemize}

\neutralbox{\textbf{Demo evidence to keep:} a valid \code{.lal} sample, a rejected \code{.txt},
all token categories, malformed input, and the generated \code{Symbol Table.txt}.
The blank rubric scores remain Ma'am's assessment.}

\clearpage
\section{3. Syntax forms for checking}
\readybox{\good\ \textbf{The EBNF and syntax checker cover the listed forms.}
The four switch rules are R42--R45 in Chapter 3.}

{\small\renewcommand{\arraystretch}{1.25}
\begin{tabularx}{\linewidth}{@{}L{29mm}Y L{35mm}@{}}
\toprule
\textbf{Family} & \textbf{Coverage Ma'am expects} & \textbf{Local status} \\\midrule
Input /1 & One input construct; its rule design is up to us. & \good\ Input rule \\
Output /3 & String only; variable value only; string plus variable value. & \good\ Three forms \\
Assignment /3 & Variable to variable; constant to variable; expression result to variable. & \good\ Three forms \\
Selection & The three IF forms and \code{SWITCH} with labeled cases and a default. & \good\ IF and SWITCH \\
Iteration /2 & Single loop and nested loops, using one chosen loop kind. & \good\ WHILE, both forms \\
Declaration & Required for us because it is already proposed. Discuss single-variable, list, and initialized forms. & \good\ Three forms \\\bottomrule
\end{tabularx}}

\profquote{One example for every rule in that statement.}{Recording 02, 01:02--01:08. Each rule needs its own matching example.}

\subsection{What each applicable rule needs for 12 November}
\begin{enumerate}
\item \textcolor{focus}{\textbf{Formal grammar components:}} terminals, nonterminals, productions,
and a grammar start symbol. The lexer's \code{q0} is not that start symbol.
\item \textcolor{focus}{\textbf{A distinct matching example:}} three examples of a single IF do not cover
three different conditional forms. Cover additional applicable rules too.
\item \textcolor{focus}{\textbf{Leftmost and rightmost derivations:}} apply the actual written productions
until the example is obtained. Show both complete derivations and one shared parse tree.
\end{enumerate}

\readybox{\good\ \textbf{Documented groundwork only:} token vocabulary, expression examples,
operator meanings, and intended precedence/associativity already exist in the specification.
Reuse these decisions, then encode them in the grammar. They are not implemented evaluation.}

\neutralbox{\textbf{Keep these distinctions clear.} \code{a = b + 3} is plain assignment;
\code{INTEGER a = b + 3} is declaration with initialization. The declaration count is not
fixed on the screen; Ma'am tentatively discusses three forms. Her printf examples do not
force C-style formatting in LAlgo.}

\clearpage
\section{4. Demonstrate the syntax checker}
\readybox{\good\ \textbf{Available foundation:} the lexer stores token type, lexeme, line, and column
in memory and creates EOF. Ma'am permits an \textbf{in-memory handoff}, one grammar terminal
at a time. Continue generating \code{Symbol Table.txt} for the lexical rubric.}

\subsection{What to show during checking}
\begin{description}[style=nextline,leftmargin=0pt,labelindent=0pt,labelwidth=0pt,labelsep=0pt,itemsep=5pt]
\item[\textcolor{ready}{1. Scan the switch sample.}]
Run \code{./build/lalgo syntax\_proposal/switch\_sample.lal}. Show the four keyword
tokens, colon delimiters, and the generated symbol table.
\item[\textcolor{ready}{2. Check the syntax.}]
Run \code{make syntax}. The checker reads the C lexer terminal stream and checks the
same EBNF printed in the proposal. It accepts nested switches and requires a final default.
\item[\textcolor{ready}{3. Show the derivations and a diagnostic.}]
Use R42--R45 for the derivations and parse trees. Remove a case colon or arm terminator
in a scratch copy to show a syntax error with its line, column, and expected token.
\item[\textcolor{focus}{4. Implement panic recovery.}]
The checker currently reports the farthest mismatch. The proposed recovery stage will
resume checking later statements after an error.
\end{description}

\profquote{The error message must be brief and concise, but it will further give the programmer the reason why there is such an error.}{Recording 04, 07:38--07:48. Include a useful reason, not just ``syntax error.''}

\subsection{Optional artifacts --- not automatic blockers}
\neutralbox{\textbf{Syntax transition diagram and parsing table are optional} in the checking guide.
Their usefulness depends on the approach we choose. Existing lexical DFA diagrams and the
\code{DELTA} transition table are different artifacts.}

\subsection{Team handoff}
\begin{enumerate}
\item \textbf{Bring the current proposal.} Chapter 3 contains the grammar and the 45
examples, each with two derivations and one shared parse tree.
\item \textbf{Show the checking evidence.} Run the lexer and syntax checker on the switch
sample. Use a malformed scratch copy for diagnostics.
\item \textbf{Prepare the next stage.} Define selector and case types, check duplicate case
labels, and add recovery that resumes after syntax errors.
\end{enumerate}

{\footnotesize\color{muted}
\textbf{Where this comes from.} The 5 November recordings 01--04, the lexical, syntax, and semantic checking guides,
and Adriel's reported iteration scope limit. Current status reflects the
switch and case revision. Quoted excerpts are lightly edited machine-assisted transcriptions; recording numbers follow
the supplied files 1--4, with timestamps restarting per file. The iteration scope limit was
reported by Adriel without a verified audio timestamp. Full transcripts and audit notes remain in the
\href{https://github.com/adr1el-m/Lalgo/blob/main/docs/meetings/2026-11-05-meeting.md}{unified meeting notes}.
Course scores are reported separately from local validation.}

\clearpage
\section{5. Ma'am's syntax checking table}
These are the criteria shown on screen. The \textbf{Existing} and \textbf{Point} columns
are left blank, just as in the guide. Our progress assessment is on pages 2--4.

{\small\setlength{\tabcolsep}{7pt}\renewcommand{\arraystretch}{1.6}
\begin{tabularx}{\linewidth}{|Y|L{22mm}|L{17mm}|}
\hline
\rowcolor{cream}\textbf{Criteria} & \textbf{Existing} & \textbf{Point} \\ \hline
1. Input must be read one by one from the symbol table (one token at a time) & & \\ \hline
2. A parser algorithm was implemented\newline Algorithm: Recursive Descent Parsing & & \\ \hline
3. Transition Diagram & & \\ \hline
4. Parsing table & & \\ \hline
5. Input Statement \quad /1 & & \\ \hline
6. Output Statement \quad /3 & & \\ \hline
7. Assignment Statement \quad /3 & & \\ \hline
8. Condition Statement \quad /3 & & \\ \hline
9. Iterative Statement \quad /2 & & \\ \hline
10. Declaration Statement & & \\ \hline
11. Error Recovery\newline Method used: & & \\ \hline
12. Error Messages & & \\ \hline
Total: & & \\ \hline
Percentage (total score/total items) & & \\ \hline
\end{tabularx}}

\neutralbox{\textbf{Read this with recording 04.} Recursive descent is an example;
items 3--4 are optional. Declaration applies when it is in the proposal, as it is for LAlgo.
Iteration /2 means single and nested loops of one chosen kind. In-memory terminal input is allowed.}

{\footnotesize\color{muted}Source: Ma'am's Syntax Analyzer screenshot, also explained in
recording 04. The written table is reproduced separately from our implementation status;
no points or percentage have been assigned here.}

\clearpage
\section{6. Ma'am's semantic checking table}
\focusbox{\todo\ \textbf{Semantic analysis still needs to be built.} We have token recognition
and source positions, but no semantic-rule checker or semantic symbol-table updates yet.}

{\small\setlength{\tabcolsep}{7pt}\renewcommand{\arraystretch}{1.25}
\begin{tabularx}{\linewidth}{|Y|L{22mm}|L{17mm}|}
\hline
\rowcolor{cream}\textbf{Criteria} & \textbf{Existing} & \textbf{Point} \\ \hline
1. Implementation of Semantic Rules:\newline
Evidence: must have a semantic rule for every syntactic rule & & \\ \hline
\multicolumn{3}{|l|}{\textbf{Semantic Errors: (Basic)}} \\ \hline
2. Input Statement\newline Data Type Mismatch\newline Variable Not Declared & & \\ \hline
3. Output Statement\newline Data Type Mismatch\newline Variable Not Declared & & \\ \hline
4. Assignment Statement\newline Data Type Mismatch\newline Variable Not Declared & & \\ \hline
5. Condition Statement\newline Data Type Mismatch\newline Variable Not Declared & & \\ \hline
6. Iterative Statement\newline Data Type Mismatch\newline Variable Not Declared & & \\ \hline
7. Declaration Statement\newline Data Type Mismatch\newline Variable Not Declared\newline
Double Declaration\newline Warning: Variable declared but not used & & \\ \hline
8. Additional Semantic Error(s) (Optional) & & \\ \hline
9. Error Recovery & & \\ \hline
10. Error Messages & & \\ \hline
11. Changes in Symbol Table value & & \\ \hline
Total Grade & & \\ \hline
\end{tabularx}}

\subsection{What we need to add for the semantic milestone}
\begin{itemize}
\item \textcolor{focus}{\textbf{A semantic rule for every syntax rule.}} Define legal types,
declarations, identifier use, and how each construct changes semantic state.
\item \textcolor{focus}{\textbf{Checks and useful diagnostics.}} Cover the listed type/declaration errors,
duplicate declarations, unused-variable warnings, and recovery. Extra error categories are optional.
\item \textcolor{focus}{\textbf{A semantic symbol table with updates.}} The current token-occurrence listing
is groundwork, not this deliverable. Define and demonstrate the value/state changes required by the guide.
\end{itemize}

{\footnotesize\color{muted}Source: the Semantic Analyzer screenshot supplied by Adriel after the
recordings. No accompanying semantic explanation/audio was supplied. The table is transcribed
as shown; exact type rules, update behavior, and scoring must be defined or clarified.
Existing/Point cells remain blank.}

\end{document}
